% ========================================
% CHƯƠNG 2: THIẾT KẾ HỆ THỐNG
% ========================================
\chapter{THIẾT KẾ HỆ THỐNG}
\label{ch:thietke}

% ----- TÓM TẮT CHƯƠNG -----
\vspace{0.5cm}
\noindent\textit{Chương này trình bày quá trình thiết kế hệ thống
Smart Home Platform, bao gồm phân tích yêu cầu chức năng và phi
chức năng, thiết kế vai trò thiết bị trong mạng Zigbee, data model
và MQTT message format, logic flow và command lifecycle, kiến trúc
cloud backend, và các quyết định trade-off quan trọng.}
\vspace{0.5cm}

% ========================================
\section{Yêu cầu chức năng và phi chức năng}
\label{sec:ch2-requirements}

\subsection{Yêu cầu chức năng}
\label{sec:ch2-func-req}

Hệ thống cần quản lý danh sách home, room và device qua REST API,
cho phép điều khiển bật/tắt đèn thông qua cluster 0x0006 từ cloud
xuống thiết bị Zigbee. Telemetry và event từ thiết bị phải được nhận
và lưu bền vững vào PostgreSQL. Vòng đời lệnh được theo dõi chặt
chẽ với các trạng thái chuẩn: accepted $\rightarrow$ queued
$\rightarrow$ sent $\rightarrow$ executed $|$ failed $|$ timeout.
Hệ thống cần phát hiện presence qua Occupancy Sensor, trong đó local
binding có sẵn nhờ Zigbee, đồng thời cho phép join thiết bị mới qua
cơ chế permit-join của Z3Gateway.

\subsection{Yêu cầu phi chức năng}
\label{sec:ch2-nonfunc-req}

Hệ thống phải deploy được trên EC2 bằng docker-compose
(\texttt{deploy/docker-compose.prod.yml}). Timeout mặc định cho
command là 5000~ms (\texttt{SB\_COMMAND\_TIMEOUT\_MS}). Broker cấu
hình ACL tách biệt principal gateway/client/monitor/bridge để giới
hạn quyền theo nguyên tắc least privilege. Hệ thống chịu được
disconnect internet ngắn hạn nhờ local control qua binding, đảm bảo
thiết bị vẫn hoạt động ngay cả khi mất kết nối cloud.

% ========================================
\section{Network và Device Roles (Zigbee)}
\label{sec:ch2-device-roles}

Coordinator/NCP chạy trên EFR32MG12 với firmware
\texttt{end\_devices/NCP}, giao tiếp host qua EZSP/ASH. Router gồm
Z3Light và Z3Switch, luôn online và relay mesh. End Device là
Occupancy Sensor, hoạt động ngắt quãng và gửi event khi phát hiện
chuyển động.

\begin{table}[H]
\centering
\caption{Ánh xạ device type và cluster chính}
\label{tab:ch2-device-roles}
\begin{tabular}{|l|l|l|l|}
\hline
\textbf{Device} & \textbf{Zigbee Role} & \textbf{Server Clusters}
  & \textbf{Client Clusters} \\ \hline
Coordinator/NCP & ZC / Trust Center & -- & On/Off Client \\ \hline
Light Device & ZR (Router) & On/Off (0x0006) & -- \\ \hline
Switch Device & ZR (Router) & -- & On/Off (0x0006) \\ \hline
Occupancy Sensor & ZED (SED) & Occupancy Sensing (0x0406) & -- \\ \hline
\end{tabular}
\end{table}

Bảng ánh xạ trên tương ứng với mô tả trong
\texttt{docs/DEVICE\_CAPABILITY\_MATRIX.md} và
\texttt{docs/ADAPTER\_ACTION\_MAP.md}. Light sử dụng cluster On/Off
0x0006 phía server, Switch sử dụng cluster On/Off phía client kết
hợp binding, còn Occupancy Sensor sử dụng cluster Occupancy Sensing.

% ========================================
\section{Data và Message Design}
\label{sec:ch2-data-design}

\subsection{MQTT Message Channels}
\label{sec:ch2-channels}

Hệ thống phân tách message thành năm kênh chính theo hướng truyền
và ngữ nghĩa.

\begin{table}[H]
\centering
\caption{Thiết kế các kênh MQTT}
\label{tab:ch2-channels}
\begin{tabular}{|l|l|l|l|}
\hline
\textbf{Channel} & \textbf{Hướng} & \textbf{Publisher}
  & \textbf{Subscriber} \\ \hline
reported & Uplink & gateway & cloud \\ \hline
event & Uplink & gateway & cloud \\ \hline
reply & Uplink & gateway & cloud \\ \hline
desired & Downlink & cloud & gateway \\ \hline
command & Downlink & cloud & gateway \\ \hline
\end{tabular}
\end{table}

\subsection{State Message Design}
\label{sec:ch2-telemetry}

Trạng thái thiết bị publish trên kênh reported mang payload chứa
\texttt{device\_id} logic, eui64 phần cứng và các attribute (ví dụ
\texttt{on\_off: true/false} cho đèn). Cloud ghi vào bảng
\texttt{device\_states} trong PostgreSQL.

\begin{lstlisting}[language={}, caption={Light reported state message}]
{
  "schema": "sb.v1",
  "msg_id": "550e8400-e29b-41d4-a716",
  "ts": "2026-01-15T10:30:00Z",
  "source": "gateway",
  "payload": {
    "device_id": "light_01",
    "device_type": "light",
    "eui64": "00:11:22:33:44:55:66:77",
    "on_off": true,
    "endpoint": 1
  }
}
\end{lstlisting}

\subsection{Command Message Design}
\label{sec:ch2-commands}

Command từ cloud mang envelope chuẩn với payload chứa action cụ thể
(ví dụ \texttt{\{"action": "on\_off", "value": true\}}).
\texttt{msg\_id} là UUID để dedup; \texttt{correlation\_id} sẽ được
gateway echo lại trong reply để cloud map về
\texttt{commands.id}.

\begin{lstlisting}[language={}, caption={Command request message}]
{
  "schema": "sb.v1",
  "msg_id": "cmd-uuid-001",
  "ts": "2026-01-15T10:31:00Z",
  "source": "cloud",
  "payload": {
    "command_id": "cmd-uuid-001",
    "target_device": "light_01",
    "action": "on_off",
    "value": true
  }
}
\end{lstlisting}

\begin{lstlisting}[language={}, caption={Command reply message}]
{
  "schema": "sb.v1",
  "msg_id": "reply-uuid-001",
  "ts": "2026-01-15T10:31:02Z",
  "source": "gateway",
  "correlation_id": "cmd-uuid-001",
  "payload": {
    "command_id": "cmd-uuid-001",
    "status": "executed",
    "message": "ZCL On/Off command sent successfully"
  }
}
\end{lstlisting}

\subsection{Versioning Rule}
\label{sec:ch2-versioning}

Namespace chứa phiên bản: \texttt{sb/v1/...}; envelope có trường
\texttt{schema: "sb.v1"}. Thay đổi không tương thích đòi hỏi bump
lên v2, giữ song song hai phiên bản trong giai đoạn chuyển tiếp.
Thêm field mới vào payload không phá vỡ consumer cũ (consumer bỏ qua
field không biết), nhưng xóa field hoặc đổi kiểu dữ liệu yêu cầu
nâng phiên bản.

% ========================================
\section{Logic Flow và State Machine}
\label{sec:ch2-logic-flow}

\subsection{Light Device State Machine}
\label{sec:ch2-light-fsm}

Thiết bị đèn chỉ có hai trạng thái kiên định ON và OFF đối với
cluster 0x0006. Trạng thái thay đổi qua ZCL command (từ gateway hoặc
switch bound trực tiếp) và phản hồi qua attribute report về
\texttt{telemetry\_rx.c}. Sau mỗi transition, Light Device cập nhật
thuộc tính \texttt{OnOff} trong On/Off Server Cluster, đổi trạng thái
LED vật lý và gửi Report Attributes lên Coordinator.

\subsection{Command Lifecycle}
\label{sec:ch2-gateway-flow}

Mỗi command đi qua một vòng đời với các trạng thái theo dõi rõ ràng:
accepted $\rightarrow$ queued $\rightarrow$ sent $\rightarrow$
executed $|$ failed $|$ timeout.

Gateway phát các chuyển tiếp accepted, queued và sent. Gateway hoặc
cloud timeout worker phát timeout. Device reply (qua gateway) phát
executed hoặc failed. Mọi reply đều mang
\texttt{correlation\_id = \{command\_id\}}, cho phép cloud cập nhật
trạng thái command trong database khi nhận reply.

\subsection{Gateway Runtime Flow}
\label{sec:ch2-app-flow}

Tiến trình Z3Gateway khởi tạo theo thứ tự: \texttt{net\_mgr} form
network, \texttt{app\_mqtt} connect broker và subscribe
desired/command, \texttt{device\_registry} khởi tạo bảng trống,
\texttt{cmd\_handler} bắt đầu loop dispatch. Khi có message command,
\texttt{cmd\_handler.c} giải mã envelope, kiểm tra schema, gọi
\texttt{device\_dispatch.c} để chọn module (\texttt{light\_ctrl} hay
\texttt{switch\_logic}), rồi gateway emit accepted/queued/sent và
chờ reply từ NCP trước khi emit executed hoặc failed.

% ========================================
\section{Cloud Backend Architecture}
\label{sec:ch2-cloud}

\subsection{Kiến trúc tổng quan}
\label{sec:ch2-cloud-arch}

Ứng dụng FastAPI được tổ chức theo layer. \texttt{main.py} khởi tạo
app cùng lifespan (database, MQTT, worker). \texttt{config.py} sử
dụng pydantic-settings với prefix \texttt{SB\_}. \texttt{database.py}
tạo async engine và session factory. \texttt{models.py} định nghĩa
bảy bảng ORM. \texttt{schemas.py} chứa Pydantic request/response.
\texttt{mqtt\_client.py} subscribe reported/event/reply rồi ghi vào
database, đồng thời expose hàm publish cho router commands.
\texttt{command\_timeout.py} là worker sweep commands quá hạn sang
trạng thái timeout. Thư mục \texttt{routers/} chứa các module health,
devices, commands và events.

\subsection{Database Design}
\label{sec:ch2-db-schema}

PostgreSQL schema (file \texttt{database/schema.sql}, tương ứng
\texttt{cloud/app/models.py}) gồm bảy bảng chính.

\begin{table}[H]
\centering
\caption{Các bảng database chính}
\label{tab:ch2-db-tables}
\begin{tabular}{|l|p{9cm}|}
\hline
\textbf{Bảng} & \textbf{Vai trò} \\ \hline
\texttt{homes} & Đơn vị tổ chức cấp cao nhất \\ \hline
\texttt{rooms} & Thuộc home, nhóm thiết bị \\ \hline
\texttt{users} & Người dùng \\ \hline
\texttt{devices} & Thiết bị Zigbee (khóa logic device\_id, kèm eui64) \\ \hline
\texttt{device\_states} & Trạng thái gần nhất của từng device \\ \hline
\texttt{events} & Log event lâu bền \\ \hline
\texttt{commands} & Vòng đời lệnh với state enum \\ \hline
\end{tabular}
\end{table}

\subsection{REST API Endpoints}
\label{sec:ch2-api}

\begin{table}[H]
\centering
\caption{REST API endpoints chính}
\label{tab:ch2-api-endpoints}
\begin{tabular}{|l|l|p{6cm}|}
\hline
\textbf{Method} & \textbf{Path} & \textbf{Ý nghĩa} \\ \hline
GET & \texttt{/health} & Healthcheck \\ \hline
GET & \texttt{/api/devices} & Liệt kê thiết bị, filter \texttt{?room\_id=} \\ \hline
GET & \texttt{/api/devices/\{device\_id\}} & Chi tiết thiết bị \\ \hline
GET & \texttt{/api/devices/\{device\_id\}/state} & Trạng thái gần nhất \\ \hline
POST & \texttt{/api/devices/\{device\_id\}/command} & Tạo lệnh, publish MQTT \\ \hline
GET & \texttt{/api/commands/\{command\_id\}} & Tra cứu lệnh theo id \\ \hline
GET & \texttt{/api/events} & Lọc event theo device\_id, event\_type, limit, offset \\ \hline
\end{tabular}
\end{table}

\subsection{MQTT -- Database Integration}
\label{sec:ch2-mqtt-db}

\texttt{mqtt\_client.py} là lớp kết nối giữa hai thế giới: subscribe
các topic uplink rồi ghi database qua session async, đồng thời expose
hàm publish để \texttt{routers/commands.py} gọi khi tạo lệnh mới.
Worker \texttt{command\_timeout.py} chạy nền, định kỳ quét bảng
commands để chuyển những bản ghi quá
\texttt{SB\_COMMAND\_TIMEOUT\_MS} sang state timeout.

% ========================================
\section{Trade-offs và Quyết định thiết kế}
\label{sec:ch2-tradeoffs}

\textbf{Loại bỏ cầu nối Python $\leftrightarrow$ gateway (tháng
4/2026).} Trước đây có một process Python
(\texttt{gateway/src/bridge.py}) làm trung gian
MQTT$\leftrightarrow$IPC qua Unix socket và NDJSON. Quyết định gộp
client MQTT vào Z3Gateway C loại bỏ một process, một format wire và
một lớp lỗi (latency IPC, reconnect hai phía). Đổi lại, mọi thay
đổi contract MQTT phải đồng bộ trong C code thay vì Python thuần.

\textbf{device\_id logic thay vì eui64 hay nwk\_addr.} eui64 là danh
tính bền nhưng gắn chặt phần cứng, nwk\_addr thì đổi mỗi lần
rejoin. Chọn \texttt{device\_id} tự sinh giúp database ổn định qua
các thay đổi hardware và tránh leak thông tin MAC.

\textbf{Retain cho state, không retain cho command.} State có thể an
toàn retain để client mới đọc được trạng thái gần nhất; command
retain sẽ gây replay sai khi gateway reconnect.

\textbf{Timeout ở cả gateway lẫn cloud.} Gateway phát timeout khi
NCP không reply; cloud worker cũng có vòng sweep phòng trường hợp
gateway không gửi được timeout (mất kết nối). Hai lớp độc lập giảm
rủi ro lệnh treo vĩnh viễn.


\newpage
